\section{\label{sec:Datasets} CT18 全球 QCD 分析概览
}

\subsection{\label{sec:summary} 概要总结}

\subsubsection{输入实验数据与最终 PDF 组合}

CT18 分析更新了广泛使用的 CT14 PDF 组~\cite{Dulat:2015mca}，对扩充后的实验测量集合（包括来自 $ep$ 对撞机 HERA 与 LHC 的高统计量数据）进行了 NNLO 与 NLO 全球拟合。CT18 实验数据集包含 ATLAS、CMS 与 LHCb 关于单举喷注、$W/Z$ 玻色子与顶夸克对产生的高统计量测量，同时保留了关键的{\it 传承}数据，如 HERA Run I 与 Run II 合并数据以及 Tevatron 的测量。
到 2018 年，LHC 各合作组已发表约三十多个可能约束 CT PDF 的实验数据集。我们用 Secs.~\ref{sec:ExptSelection} 与 \ref{sec:Advances} 综述的方法挑选了截至 2018 年年中可获得的最有前景的实验，随后在完整拟合框架内深入考察了这些数据集的影响。Sec.~\ref{sec:data_overview} 给出这些实验的概述。
CT18 所含数据点的运动学分布作为典型部分子动量分数 $x$ 与 QCD 因子化标度 $\Q$ 的函数示于图~\ref{fig:ct18data_xmu}。\footnote{典型动量分数与因子化标度的估计方法见文献~\cite{Wang:2018heo}。}与全球 PDF 拟合长期以来的情形一样，纳入的数据覆盖了 $x$--$\Q$ 平面上很大的运动学范围。
\begin{figure}[tb]
	\includegraphics[width=0.99\textwidth]{images/ct18data_xmu.pdf}
	\caption{CT18 数据集，表示在部分子 $(x,\Q)$ 平面中，基于 Born 级运动学匹配 $(x,\Q) = (x_B, Q)$（DIS 等）。此处采用的匹配约定见文献~\cite{Wang:2018heo}。图中还标出了在 CT18Z 中拟合的 ATLAS 7 TeV $W/Z$ 产生数据（ID=248），记作 ATL7WZ'12。
		\label{fig:ct18data_xmu}}
\end{figure}

鉴于某些测量达到前所未有的精度，最新的 LHC 数据必须用微扰 QCD 的 NNLO 理论预言来分析。
本分析得到的拟合 PDF 绘于图~\ref{fig:ct18pdf}：上面板为两个相距甚远的标度 $\Q\!=\!2$ 与 $\,100$ GeV 处（分别为左右两图）的 CT18 PDF；下面板为我们合并的替代分析 CT18Z 给出的相应 PDF。

\begin{figure}[htbp]
\hspace*{-0.9cm}\includegraphics[width=0.52\textwidth]{images/pdfs_CT18___2GeV_Epdfs_cus-lin_ect.pdf}
\includegraphics[width=0.52\textwidth]{images/pdfs_CT18_100GeV_Epdfs_cus-lin_ect.pdf}
\hspace*{-0.9cm}\includegraphics[width=0.52\textwidth]{images/pdfs_CT18ZNNLO___2GeV_A90CL_Epdfs_cus-lin-d5_ect.pdf}
\includegraphics[width=0.52\textwidth]{images/pdfs_CT18ZNNLO_100GeV_A90CL_Epdfs_cus-lin-d5_ect.pdf}
\caption{上面板：$\Q\! =\! 2$ GeV 与 $\Q\! =\! 100$ GeV 处 $u, \overline{u}, d, \overline{d}, s = \overline{s}$ 与 $g$ 的 CT18 部分子分布函数。
		下面板：CT18Z 的类似曲线。所有情形中胶子 PDF 都按 $g(x,\Q)/5$ 缩小显示。粲分布 $c(x,\Q)$ 由演化微扰生成（CT18 与 CT18Z 的初始标度分别为 $Q_0\! =\! 1.3$ 与 1.4 GeV），图中一并给出。
		\label{fig:ct18pdf}}
\end{figure}

最终的 CT18(Z) 数据组合共包含 $N_\mathit{pt}\! =\!$ 3681（3493）个数据点，NNLO 下给出 $\chi^2/N_\mathit{pt}=1.17\,(1.19)$。
PDF 不确定度按照 CT14 全球分析~\cite{Dulat:2015mca} 的两级判据在 90\% 概率水平构造。这些 PDF 在取世界平均值 $\alpha_s(M_Z)=0.118$~\cite{Tanabashi:2018oca} 的假定下获得。组合的 PDF+$\alpha_s$ 不确定度可用每族 PDF 的专门 $\alpha_s$ 系列、把 PDF 与 $\alpha_s$ 不确定度按平方和叠加来计算，说明见文献~\cite{Lai:2010nw}。

在四个 PDF 组合（CT18、A、X、Z）之中，CT18 与 CT18Z 在 PDF 形状上差异最大，尤其是拟合的胶子与奇异性分布 $g(x,Q)$、$s(x,\Q)$ 的 $x$ 依赖性以及部分 PDF 不确定度。随着下文讨论的 LHC Run-1 数据的纳入，CT18 的胶子密度 $g(x,\Q)$ 精度较 CT14 有适度改进。而对 CT18Z，由于 Sec.~\ref{sec:alt} 与 App.~\ref{sec:AppendixCT18Z} 描述的对 DIS 数据处理的修改，胶子及微扰生成的粲 PDF 的不确定度有所增大，特别是在最低的 $x\! <\! 10^{-3}$ 区域。
这些最终 PDF 依赖实验数据的众多系统性因素。为得到可信的 PDF 不确定度估计，必须细致考察这些系统效应；数值计算的规模也需要扩大。

\subsubsection{\label{sec:summary-HERA2} HERA I+II 合并 DIS 数据与依赖 $x_B$ 的因子化标度}
即便在 LHC 时代，来自 $ep$ 对撞机 HERA 的 DIS 数据仍为 CT18 PDF 提供最主要的约束。这一主导地位由 \texttt{ePump}、\texttt{PDFSense} 与拉格朗日乘子方法的独立应用共同揭示。
CT18 采用 HERA Run-I 与 Run-II DIS 的最终（``合并''）数据集~\cite{Abramowicz:2015mha}，它取代了 CT14~\cite{Dulat:2015mca} 所用的仅含 HERA Run-I 的数据集~\cite{Aaron:2009aa}。
基于对最终 HERA 数据的拟合，曾发布过渡性 PDF 组 \CTHERAII~\cite{Hou:2016nqm}。我们发现 HERA I+II 数据与 CT14 及 \CTHERAII{} PDF 总体符合良好，且两个 PDF 组合都能同样好地描述全球分析中的非 HERA 数据。
与此同时，我们观察到 HERA I+II 数据集中 $e^{+}p$ 与 $e^{-}p$ DIS 截面之间存在一定分歧（``统计张力''）。我们的结论是，目前无法对这些张力在全 Bjorken $x$ 与 HERA 轻子-质子动量转移 $Q$ 可及范围内的分布模式给出合理解释。
利用 CT18 拟合扩展这些研究时，我们考察了 QCD 标度选择对小 $x_B$ 区单举 DIS 数据的影响，详见后文 Sec.~\ref{sec:alt}。

\begin{figure}[tb]
	\includegraphics[width=0.59\textwidth]{images/ct18saturationToct18.pdf}
	\includegraphics[width=0.4\textwidth]{images/HERAchi2NptWT_scan2Tct18z.pdf}
\caption{左：DIS 数据集分别用依赖 $x_B$ 的因子化标度与标准因子化标度拟合所得候补 CT18 NNLO PDF 的比值。右：采用依赖 $x_B$ 的 DIS 因子化标度的 CT18Z 拟合中四个 HERA 数据集的 $\chi^2/N_{pt}$ 值随 HERA I+II 单举 DIS 数据集统计权重的变化。}
\label{fig:saturation}
\end{figure}

我们发现，用一个特殊的、依赖于 Bjorken $x_B$（而非动量分数 $x$）的因子化标度 $\mu_{F,x}$ 计算 DIS 的 NNLO 理论截面，可使 HERA 数据的拟合质量改善约 50 个单位；该标度将在 Sec.~\ref{sec:alt} 中引入。图~\ref{fig:saturation}（左）展示了用因子化标度 $\mu_{F,x}$ 拟合 DIS 数据集所得候补 CT18 PDF 相对于采用名义标度 $\mu_F=Q$ 的 CT18 PDF 的变化。采用标度 $\mu_{F,x}$ 后，与名义 CT18 拟合相比，我们观察到 $x < 10^{-2}$ 处 $u$、$d$（反）夸克 PDF 减小而胶子与奇异性 PDF 增大；同时在无约束区域 $x > 0.5$ 同一些 PDF 出现若干补偿性变化，以满足价夸克与动量求和规则。

图~\ref{fig:saturation} 右面板给出四个 HERA 数据集（单举中性流与带电流 DIS~\cite{Abramowicz:2015mha}、约化粲与底产生截面，以及 H1 纵向结构函数 $F_L(x_B,Q^2)$~\cite{Collaboration:2010ry}）在拟合中的 $\chi^2/N_{pt}$ 值（$\chi^2$ 除以数据点数 $N_\mathit{pt}$）随 HERA I+II 单举 DIS 数据集统计权重 $w$ 的变化~\cite{Abramowicz:2015mha}。默认 CT18Z 拟合对应 $w=1$；当 $w=10$ 时，CT18Z 拟合越来越表现为仅含 HERA 的拟合。可以看到，取标度 $\mu^2_{F,x}$ 且 $w=10$ 时，单举 DIS 数据集的 $\chi^2/N_{pt}$ 几乎改善到``重求和''型仅-HERA 拟合（无内禀粲）所见的水平~\cite{Ball:2017otu,Abdolmaleki:2018jln}。对粲半单举 DIS（SIDIS）截面与 H1 $F_L$ 的拟合质量也得到改善。\footnote{同时使用单独的 H1 $F_L$ 数据与 HERA-II 合并数据会引入一定的重复计数。但我们已核查这一选择不会明显改变 PDF，却能为小 $x$ 区拟合优度提供有用的指标。特别值得注意的是，与 CT18 相比，H1 $F_L$ 数据（ID=169）的总 $\chi^2$ 值在 CT18X 与 CT18Z 拟合中变小而非变大。}

H1 与 ZEUS 合作组在文献~\cite{H1:2018flt}（2018 年）发表的新合并粲、底产生测量已被研究；就当前版本而言，若以这些测量替换 CT18 全球分析中此前的测量，则无法以合理的 $\chi^2$ 拟合之。
而且，这些新合并数据与若干 CT18 数据集之间观察到轻度张力，如 LHCb 7 与 8 TeV $W/Z$ 产生数据、CDF run II 的 $Z$ 快度数据、CMS 8 TeV 单举喷注产生，以及 $t\bar{t}$ 双微分 $p_T$ 与 $y$ 截面。因此我们决定不把这些数据纳入 CT18 全球分析，它们需要专门研究。
在文献~\cite{H1:2018flt} 的 H1+ZEUS 分析中，这些测量的 $\chi^2$ 也并非最优。这被归因于中间/小 $x$ 区数据与理论斜率的差异。在尝试拟合新合并粲、底产生测量时，我们注意到对中间/小 $x$ 处更硬胶子的偏好。
我们目前正在单独研究这些数据，尤其正在探索新关联系统不确定度的影响——它由旧版数据的 42 个增加到新版的 167 个。
这项新研究的结果将在另一篇即将发表的论文中给出。


\begin{figure}[tb]
	\includegraphics[width=0.6\textwidth]{images/sqrt2chi2_ct18nnloCount.pdf}
	\caption{遍布全部 CT18 数据集的有效高斯变量（$S_E$）的直方图。
	两个方块与两颗星分别标记 NuTeV 双缪子与 CCFR 双缪子数据的 $S_E$ 值。
\label{fig:sn_ct18}}
\end{figure}

\subsubsection{遴选新的 LHC 实验 \label{sec:ExptSelection} }
为 CT18 拟合挑选最有前景的 LHC 实验时，我们必须应对一个反复出现的难题——来自 HERA、LHC 与 Tevatron 的最新实验数据的各（子）集合之间存在的统计张力。对撞机数据迅速改善的精度揭示出以往无关紧要的实验或理论反常。施加强拟合优度检验即可显现这些反常~\cite{Kovarik:2019xvh}。图~\ref{fig:sn_ct18} 用基于有效高斯变量 $S_E\equiv \sqrt{2\chi_E^2}-\sqrt{2 N_E-1}$~\cite{Lai:2010vv} 的表示展示张力的程度；该变量由单个数据集 $E$ 的 $\chi^2$ 值与数据点数 $N_\mathit{E}$ 构造。在一个理论与数据之差符合高斯随机涨落的理想拟合中，$S_E$ 的概率分布应近似标准正态分布（半宽为 1）。而在 CTEQ-TEA 及外部组的全球拟合中，我们看到的是如图~\ref{fig:sn_ct18} 那样更宽的 $S_E$ 分布，其中一些最全面、最精密的数据集（特别是 HERA I+II 单举 DIS~\cite{Abramowicz:2015mha} 与 ATLAS 7 TeV $W/Z$ 产生~\cite{Aaboud:2016btc}）的 $S_E$ 值高达五个单位甚至更多。于是问题在于：如何从一个不断增长的测量清单中为全球分析挑出干净而精确的实验，同时最大程度地保持所选实验之间的相容性？
例如，有许多潜在对 PDF 敏感的 LHC 数据集~\cite{Rojo:2015acz}，包括涉及高 $p_{T}$ $Z$ 玻色子、$t\bar{t}$ 对、孤立光子以及小 $x$ 重味（粲或底）夸克产生的新型测量。把所有这类候选实验都放进完整的全球拟合并不现实：在 NNLO 下 CPU 开销随实验数据集数目快速增长；拟合不佳的实验只会增大而非减小最终 PDF 不确定度。在 CT14 拟合 33 个实验时，单线程模式下生成一套误差 PDF 需要数天 CPU 时间。因此在这种配置下再增加 20--30 个实验是不可能的。
CTEQ-TEA 组通过多管齐下的努力解决了这些难题，使我们得以纳入十一个关于 $W^\pm$、$Z$、喷注与 $t\bar{t}$ 产生的 7、8 TeV 新 LHC 数据集。


\subsubsection{拟合方法学的进展 \label{sec:Advances} }
为确定适合全球拟合的候选实验数据集，我们开发了两个快速初步分析程序。\texttt{PDFSense} 程序~\cite{Wang:2018heo} 由南方卫理公会大学（SMU）开发，用于在做拟合之前定量预测哪些数据集会对全球 PDF 拟合产生影响。\texttt{ePump} 程序~\cite{Schmidt:2018hvu} 由密歇根州立大学（MSU）开发，应用 Hessian 剖析在全球拟合前快速估计数据对 PDF 的影响。这些程序完全基于先前发表的 Hessian 误差 PDF，为挑选最有价值的实验提供了有益指引。Sec.~\ref{sec:Sensitivity} 展示 \texttt{PDFSense} 的一个应用。

正如将在 Appendix~\ref{sec:Appendix4xFitter} 讨论的，常用版本 2.0.0 的 \texttt{xFitter} 程序~\cite{Alekhin:2014irh} 所内置的开箱即用 Hessian 剖析算法与 CTEQ-TEA 的 PDF 不确定度定义不一致，在多项剖析 CTEQ-TEA PDF 的研究中给出了过于乐观的 $\chi^2$ 值与 PDF 不确定度。\texttt{ePump} 程序没有这个缺陷。其 Hessian 更新算法能更好地复现完整 CT14 与 CT18 拟合中数据集的 $\chi^2$ 值以及相应的 PDF 不确定度——后者按 CTEQ-TEA 自 CT10 NLO~\cite{Gao:2013xoa} 起采用的两级 $\chi^2$ 定义给出。

CT 拟合代码已并行化，以便在高性能计算集群上更快完成拟合（一次拟合由许多天缩短到几小时）。对尽可能多的相关 LHC 数据，我们用 \texttt{APPLGrid/fastNLO} 表~\cite{Kluge:2006xs,Carli:2010rw}（再乘以制表的逐点 NNLO/NLO $K$ 因子修正）计算了多种新 LHC 过程的 NLO 截面：$W/Z$ 玻色子、高 $p_{T}$ $Z$ 玻色子与单举喷注的产生；对 LHC 的 $t\bar{t}$ 对产生则用 \texttt{fastNNLO} 表~\cite{Czakon:2017dip,fastnnlo:grids} 计算 NNLO 截面。\texttt{APPLgrid} 表与其他组的同类公开表格做了交叉验证，并在速度与精度上作了优化。

\subsubsection{理论不确定度与参数化不确定度的估计}
我们花了大量精力理解 PDF 不确定度的来源。对标度选择相关的理论不确定度，就受影响的诸如 DIS、单举喷注与高 $p_{T}$ $Z$ 玻色子产生等过程作了研究。
考虑的其他理论不确定度包括不同 NNLO 与重求和程序之间的差异（如
\texttt{DYNNLO}~\cite{Catani:2007vq,Catani:2009sm}、
\texttt{MCFM}~\cite{Campbell:2010ff, Boughezal:2016wmq, MCFM8}、
\texttt{FEWZ}~\cite{Gavin:2010az, Gavin:2012sy, Li:2012wna}、
\texttt{NNLOJet}~\cite{Ridder:2015dxa,Gehrmann-DeRidder:2017mvr,Currie:2016bfm,Currie:2017ctp}、
\texttt{ResBos}~\cite{Ladinsky:1993zn, Balazs:1997xd}），以及蒙特卡洛（MC）积分误差，见 Sec.~\ref{sec:calcs}。
具体而言，我们在 CT18(Z) 分析中对单举喷注与高 $p_T$ $Z$ 玻色子产生数据计入了 MC 误差；但与 QCD 标度选择及理论计算程序选择相关的 PDF 不确定度并未被系统地计入本分析。
重要的 PDF 参数化不确定度通过用 $\mathcal{O}(250)$ 种试验函数形式重复拟合加以研究。[CT10 之后的拟合用 Bernstein 多项式参数化 PDF，这便于尝试大范围参数化形式来量化/消除潜在偏差。Appendix~\ref{sec:AppendixParam} 给出此类参数化的一个例子。]名义 CT18 PDF 组的最终不确定度的确定方式是覆盖备选参数化形式、备选拟合设定或标度选择下获得的中心解，见 Sec.~\ref{sec:Paramstudies}。

\subsection{CT18 拟合的实验数据集}
\label{sec:data_overview}

CT18 全球分析以 \CTHERAII~\cite{Hou:2016nqm} 的数据集基线为起点，加上 2018 年年中之前发表的 LHC 结果。\CTHERAII{} 基线中的实验列于表 \ref{tab:EXP_1}，CT18(Z) 拟合纳入的新 LHC 数据集见表 \ref{tab:EXP_2}。
表~\ref{tab:EXP_1} 与~\ref{tab:EXP_2} 还列出了全球拟合中每个数据集的数据点数、$\chi^2$ 与有效高斯变量 $S_E$。
大多数数据集被包含在全部四个 PDF 组合中；具体选择的差异将在出现处指明。

\subsubsection{勘绘新数据集对 PDF 的敏感度 \label{sec:Sensitivity} }

如 Secs.~\ref{sec:ExptSelection} 与 \ref{sec:Advances} 所述，我们采用了一种基于 Hessian 敏感度变量的新方法~\cite{Wang:2018heo,Hobbs:2019gob}（它是理论观测量与 PDF 之间 Hessian 关联~\cite{Pumplin:2001ct,Nadolsky:2001yg,Nadolsky:2008zw} 的信息化后代），定量确定数据对全球拟合及特定截面的影响层级。

\begin{figure}[tb]
	\includegraphics[width=0.48\textwidth]{images/Cf_higgs.pdf}
	\includegraphics[width=0.48\textwidth]{images/Sf_higgs.pdf}
\caption{
	左：考虑纳入 CT18 的候选数据，按 14 TeV 总希格斯截面 $\sigma_{H}(14\, \mathrm{TeV})$ 与各点残差（在 \texttt{PDFSense} 框架内确定~\cite{Wang:2018heo}）之间的 Pearson 关联 $C_f$ 的大小评估。
	右：对 CT18 候选数据的类似评估，但改用{\it 敏感度} $|S_f|(x_i,\Q_i)$ 计算。两个面板中都施加了高亮割断，以突出按各度量排位最高的约 $\sim\!300$ 个点。
  }
\label{fig:PDFSense}
\end{figure}

作为一个演示，\texttt{PDFSense} 框架可以预先预测哪些被拟合数据集对 LHC 最关键的预言之一——例如经由 $gg$ 融合过程的希格斯玻色子截面（$\sigma_H$，$\sqrt{s}=14$ TeV）——可能有最大影响。要判断哪些数据集对此类截面影响最大，通常做法是考察实验数据点与负责希格斯产生的运动学区域内胶子分布的 Pearson 关联~\cite{Pumplin:2001ct,Nadolsky:2001yg,Nadolsky:2008zw}。图~\ref{fig:PDFSense} 左侧给出与 14 TeV 希格斯玻色子截面绝对关联 $|C_f|$（按文献~\cite{Wang:2018heo} 由统计残差定义）最高的数据点。

按这一度量，可能存在多个高影响数据集， notably HERA 中性流 DIS、LHC 与 Tevatron 喷注产生，以及 HERA 粲夸克产生。然而关联并不反映数据点的实验不确定度：某个实验截面可能与负责希格斯产生的 $x$ 区间内的胶子分布高度相关，但若实验不确定度大，它对希格斯玻色子截面的约束未必强。反过来，一个关联不那么大但（统计与系统）不确定度更小的实验截面可能提供更强的约束。

因此约束水平由文献~\cite{Wang:2018heo} 定义的敏感度变量 $S_f$ 更好地预示。CT18 全球拟合所用数据点中对 14 TeV $\sigma_H$ 的 PDF 依赖具有最高绝对敏感度 $|S_f|$ 者示于图~\ref{fig:PDFSense} 右侧。具有高敏感度的数据点（来自更多数目的实验）比高关联识别出的更多。除 HERA I+II 的 DIS 数据外，还有固定靶 DIS 实验的贡献以及 LHC 的测量。

正如 Sec.~\ref{sec:LMScans} 将用拉格朗日乘子扫描所示，HERA I+II 数据集凭借其丰富的数据点与较小的实验误差，仍主导着对 LHC 希格斯产生敏感区间内胶子分布的约束。
由于旧数据集持续的影响力，我们将发现 CT18 中希格斯玻色子产生 PDF 不确定度的缩减不如 CT14 显著。此外，部分最敏感数据集之间的张力限制了希格斯截面不确定度的缩减。这些效应将在 Sec.~\ref{sec:LMScans} 详细探讨。

下面我们讨论 CT18 分析纳入的新数据集，并突出替代拟合的差异。


\subsubsection{基线数据集 \label{sec:Baseline}}

CT18 全球分析从 \CTHERAII{} 继承了一批列于表 1 的高精度非-LHC 实验。其中 HERA I+II DIS 数据集提供最显著的约束，其次是一组固定靶中性流 DIS 实验：BCDMS、NMC 与 CCFR。
类似地，此前已纳入的一批中微子 DIS 测量为海（反）夸克提供有价值的约束。
在这些实验中，我们发现由 CDHSW 中微子-铁深度非弹性散射提取的单核子结构函数 $F_2^p$ 与 $F_3^p$ 相比 CCFR 及其他实验表现出对 $x \gtrsim 0.1$ 处更硬胶子 PDF 的偏好，参见图~\ref{fig:LMg18}。
这一众所周知的行为反映 CDHSW 测得的 $F_2^p$ 与 $x F_3^p$ 具有比 CCFR 类似测量~\cite{Seligman:1997fe} 更大的对数斜率，而这又可能反映两个实验间能量刻度与分辨涂抹的差异~\cite{Bodek:1996zu}。
因此，为帮助获得近期 LHC 数据偏好的较软大 $x$ 胶子行为，我们把 CDHSW $F_2$ 与 $x_B F_3$ 数据集从 CT18Z 分析中剔除，而其余 CT18 PDF 组合仍包含它们。

我们继续纳入来自 Tevatron 与固定靶实验的各类轻子对产生测量，汇总于表~\ref{tab:EXP_1}。低统计量的 LHCb 7 TeV $W/Z$ 产生数据~\cite{Aaij:2012vn} 以及 ATLAS、CMS 7 TeV 喷注产生数据~\cite{Aad:2011fc, Chatrchyan:2012bja} 在 CT18 分析中被更新的测量取代，见下一节概述。



\begingroup

\begin{table}[htbp]
\begin{tabular}{|l|lr|c|c|c|c|}
\hline
\textbf{ID\# }  & \textbf{Experimental data set} &  & $N_{pt, E}$  & $\chi^2_E$ & $\chi^{2}_E/N_{pt, E}$  & $S_E$ \tabularnewline
\hline
\hline
 160 & HERAI+II 1 fb$^{-1}$, H1 and ZEUS NC and CC $e^\pm p$ reduced cross sec. comb.      & \cite{Abramowicz:2015mha}   & 1120  &  1408(1378) &   1.3( 1.2) &   5.7(  5.1)   \tabularnewline\hline
 101 & BCDMS $F_{2}^{p}$                                                                   & \cite{Benvenuti:1989rh}     &  337  &   374 ( 384) &   1.1( 1.1) &   1.4(  1.8)   \tabularnewline\hline
 102 & BCDMS $F_{2}^{d}$                                                                   & \cite{Benvenuti:1989fm}     &  250  &   280 ( 287) &   1.1( 1.1) &   1.3(  1.6)   \tabularnewline\hline
 104 & NMC $F_{2}^{d}/F_{2}^{p}$                                                           & \cite{Arneodo:1996qe}       &  123  &   126 ( 116) &   1.0( 0.9) &   0.2( -0.4)   \tabularnewline\hline
 108$^{\dagger}$ & CDHSW $F_{2}^{p}$                                                       & \cite{Berge:1989hr}         &   85  &    85.6 ( 86.8) &   1.0( 1.0) &   0.1(  0.2)   \tabularnewline\hline
 109$^{\dagger}$ & CDHSW $x_B F_{3}^{p}$                                                       & \cite{Berge:1989hr}         &   96  &    86.5 (  85.6) &   0.9( 0.9) &  -0.7( -0.7)   \tabularnewline\hline
 110 & CCFR $F_{2}^{p}$                                                                    & \cite{Yang:2000ju}          &   69  &    78.8(  76.0) &   1.1( 1.1) &   0.9(  0.6)   \tabularnewline\hline
 111 & CCFR $x_B F_{3}^{p}$                                                                   & \cite{Seligman:1997mc}      &   86  &    33.8(  31.4) &   0.4( 0.4) &  -5.2( -5.6)   \tabularnewline\hline
 124 & NuTeV $\nu\mu\mu$ SIDIS                                                             & \cite{Mason:2006qa}         &   38  &    18.5(  30.3) &   0.5( 0.8) &  -2.7( -0.9)   \tabularnewline\hline
 125 & NuTeV $\bar\nu \mu\mu$ SIDIS                                                        & \cite{Mason:2006qa}         &   33  &    38.5(  56.7) &   1.2( 1.7) &   0.7(  2.5)   \tabularnewline\hline
 126 & CCFR $\nu\mu\mu$ SIDIS                                                              & \cite{Goncharov:2001qe}     &   40  &    29.9(  35.0) &   0.7( 0.9) &  -1.1( -0.5)   \tabularnewline\hline
 127 & CCFR  $\bar\nu \mu\mu$ SIDIS                                                        & \cite{Goncharov:2001qe}     &   38  &    19.8(  18.7) &   0.5( 0.5) &  -2.5( -2.7)   \tabularnewline\hline
 145 & H1 $\sigma_{r}^{b}$                                                                 & \cite{Aktas:2004az}         &   10  &     6.8(   7.0) &   0.7( 0.7) &  -0.6( -0.6)   \tabularnewline\hline
 147 & Combined HERA charm production                                                      & \cite{Abramowicz:1900rp}    &   47  &    58.3(  56.4) &   1.2( 1.2) &   1.1(  1.0)   \tabularnewline\hline
 169 & H1 $F_{L}$                                                                          & \cite{Collaboration:2010ry} &    9  &    17.0(  15.4) &   1.9( 1.7) &   1.7(  1.4)   \tabularnewline\hline
 201 & E605 Drell-Yan process                                                              & \cite{Moreno:1990sf}        &  119  &   103.4( 102.4) &   0.9( 0.9) &  -1.0( -1.1)   \tabularnewline\hline
 203 & E866 Drell-Yan process  $\sigma_{pd}/(2\sigma_{pp})$                                & \cite{Towell:2001nh}        &   15  &    16.1(  17.9) &   1.1( 1.2) &   0.3(  0.6)   \tabularnewline\hline
 204 & E866 Drell-Yan process  $Q^3 d^2\sigma_{pp}/(dQ dx_F)$                              & \cite{Webb:2003ps}          &  184  &   244 ( 240) &   1.3( 1.3) &   2.9(  2.7)   \tabularnewline\hline
 225 & CDF Run-1 lepton $A_{ch}$, $p_{T\ell}>25$ GeV                                    & \cite{Abe:1998rv}           &   11  &     9.0(   9.3) &   0.8( 0.8) &  -0.3( -0.2)   \tabularnewline\hline
 227 & CDF Run-2 electron $A_{ch}$, $p_{T\ell}>25$ GeV                                     & \cite{Acosta:2005ud}        &   11  &    13.5(  13.4) &   1.2( 1.2) &   0.6(  0.6)   \tabularnewline\hline
 234 & D\O~ Run-2 muon $A_{ch}$, $p_{T\ell}>20$ GeV                                        & \cite{Abazov:2007pm}        &    9  &     9.1(   9.0) &   1.0( 1.0) &   0.2(  0.1)   \tabularnewline\hline
 260 & D\O~ Run-2 $Z$ rapidity                       &
 \cite{Abazov:2007jy}        &   28  &    16.9(  18.7) &   0.6( 0.7) &  -1.7( -1.3)   \tabularnewline\hline
 261 & CDF Run-2 $Z$ rapidity                                                              & \cite{Aaltonen:2010zza}     &   29  &    48.7(  61.1) &   1.7( 2.1) &   2.2(  3.3)   \tabularnewline\hline
 266 & CMS 7 TeV $4.7\mbox{ fb}^{-1}$, muon $A_{ch}$, $p_{T\ell}>35$ GeV                   & \cite{Chatrchyan:2013mza}   &   11  &     7.9(  12.2) &   0.7( 1.1) &  -0.6(  0.4)   \tabularnewline\hline
 267 & CMS 7 TeV $840\mbox{ pb}^{-1}$, electron $A_{ch}$, $p_{T\ell}>35$ GeV               & \cite{Chatrchyan:2012xt}    &   11  &    4.6(  5.5) &   0.4( 0.5) &   -1.6(  -1.3)   \tabularnewline\hline
 268$^{\ddagger\ddagger}$ & ATLAS 7 TeV $35\mbox{ pb}^{-1}$ $W/Z$ cross sec., $A_{ch}$                          & \cite{Aad:2011dm}           &   41  &    44.4   (50.6) &   1.1( 1.2) &   0.4(  1.1)   \tabularnewline\hline
 281 & D\O~ Run-2 $9.7 \mbox{ fb}^{-1}$ electron $A_{ch}$, $p_{T\ell}>25$ GeV              & \cite{D0:2014kma}           &   13  &    22.8(  20.5) &   1.8( 1.6) &   1.7(  1.4)   \tabularnewline\hline
 504 & CDF Run-2 inclusive jet production                                                  & \cite{Aaltonen:2008eq}      &   72  &   122 ( 117) &   1.7( 1.6) &   3.5(  3.2)   \tabularnewline\hline
 514 & D\O~ Run-2 inclusive jet production                                                 & \cite{Abazov:2008ae}        &  110  &   113.8 ( 115.2) &   1.0( 1.0) &   0.3(  0.4)   \tabularnewline
\hline
\hline
\end{tabular}
\caption{CT18(Z) NNLO 全球分析包含的数据集。这里依据 $\chi^2_E$、$\chi^2_E/N_{pt, E}$ 与 $S_E$ 直接比较 CT18 NNLO 与 CT18Z NNLO 的拟合质量；$N_{pt, E}$ 与 $\chi^2_E$ 是实验 $E$ 在全局极小处的点数与 $\chi^2$ 值。$S_E$ 是量化与各实验符合程度的有效高斯参数~\cite{Lai:2010vv, Gao:2013xoa, Dulat:2013hea}。标记 $\ddagger\ddagger$ 的 ATLAS 7 TeV $35\mbox{ pb}^{-1}$ $W/Z$ 数据集在 CT18A 与 CT18Z 拟合中被更新版本（4.6 fb$^{-1}$）替换。标记 $\dagger$ 的 CDHSW 数据未纳入 CT18Z 拟合。括号内数字为 CT18Z NNLO 拟合的结果。
\label{tab:EXP_1} }
\end{table}
\endgroup



\begingroup

\begin{table}[tb]
\begin{tabular}{|l|lr|c|c|c|c|}
\hline
\textbf{ID\# }  & \textbf{Experimental data set} &  & $N_{pt, E}$  & $\chi^2_E$ & $\chi^{2}_E/N_{pt, E}$  & $S_E$ \tabularnewline
\hline
\hline
 245 & LHCb 7 TeV 1.0 fb$^{-1}$ $W/Z$ forward rapidity cross sec.                          & \cite{Aaij:2015gna}        &  33  &   53.8 (  39.9) &  1.6 (  1.2) &  2.2 ( 0.9) \tabularnewline\hline
 246 & LHCb 8 TeV  2.0 fb$^{-1}$ $Z\rightarrow e^{-} e^{+}$ forward rapidity cross sec.    & \cite{Aaij:2015vua}        &  17  &   17.7 (  18.0) &  1.0 (  1.1) &  0.2 ( 0.3) \tabularnewline\hline
 248$^{\ddagger}$ & ATLAS 7 TeV 4.6 fb$^{-1}$, $W/Z$ combined cross sec.                   & \cite{Aaboud:2016btc}      &  34  &  287.3 (  88.7) &  8.4 (  2.6) & 13.7 ( 4.8) \tabularnewline\hline
 249 & CMS 8 TeV 18.8 fb$^{-1}$ muon charge asymmetry $A_{ch}$                                & \cite{Khachatryan:2016pev} &  11  &   11.4 (  12.1) &  1.0 (  1.1) &  0.2 ( 0.4) \tabularnewline\hline
 250 & LHCb 8 TeV 2.0fb$^{-1}$ $W/Z$ cross sec.                                            & \cite{Aaij:2015zlq}        &  34  &   73.7 (  59.4) &  2.1 (  1.7) &  3.7 ( 2.6) \tabularnewline\hline
 253 & ATLAS 8 TeV 20.3 fb$^{-1}$, $Z$ $p_T$ cross sec.                                    & \cite{Aad:2015auj}         &  27  &   30.2 (  28.3) &  1.1 (  1.0) &  0.5 ( 0.3) \tabularnewline\hline
 542 & CMS 7 TeV 5 fb$^{-1}$, single incl. jet cross sec., $R=0.7$ (extended in y)         & \cite{Chatrchyan:2014gia}  & 158  &  194.7 ( 188.6) &  1.2 (  1.2) &  2.0 ( 1.7) \tabularnewline\hline
 544 & ATLAS 7 TeV  4.5 fb$^{-1}$, single incl. jet cross sec., $R=0.6$                    & \cite{Aad:2014vwa}         & 140  &  202.7 ( 203.0) &  1.4 (  1.5) &  3.3 ( 3.4) \tabularnewline\hline
 545 & CMS 8 TeV 19.7 fb$^{-1}$, single incl. jet cross sec., $R=0.7$, (extended in y)     & \cite{Khachatryan:2016mlc} & 185  &  210.3 ( 207.6) &  1.1 (  1.1) &  1.3 ( 1.2) \tabularnewline\hline
 573 & CMS 8 TeV 19.7 fb$^{-1}$, $t\bar{t}$ norm. double-diff. top $p_T$ and $y$ cross sec. & \cite{Sirunyan:2017azo}    &  16  &   18.9 (  19.1) &  1.2 (  1.2) &  0.6 ( 0.6) \tabularnewline\hline
 580 & ATLAS 8 TeV 20.3 fb$^{-1}$, $t\bar{t}$ $p_{T}^{t}$ and $m_{t\bar{t}}$ abs. spectrum    & \cite{Aad:2015mbv}         &  15  &    9.4 (  10.7) &  0.6 (  0.7) & -1.1 (-0.8) \tabularnewline
 \hline
\end{tabular}
\caption{同表~\ref{tab:EXP_1}，针对新纳入的 LHC 测量。
	标记 $\ddagger$ 的 ATLAS 7 TeV $W/Z$ 数据（4.6 fb$^{-1}$）纳入 CT18A 与 CT18Z 全球拟合，但不纳入 CT18 与 CT18X。
\label{tab:EXP_2} }
\end{table}
\endgroup





\subsubsection{来自 $W/Z$ 矢量玻色子产生的 LHC 精密数据
\label{sec:DataWZ}
}


CT18(Z) 全球分析使用来自 LHC Run-I 的 $W/Z$ 矢量玻色子产生数据，包括 ATLAS、CMS 与 LHCb 合作组的测量。

拟合中纳入的 ATLAS 测量为：

\begin{itemize}

	\item 积分亮度 4.6 fb$^{-1}$、$\sqrt{s}=7$ TeV 的 $W/Z$ 合并截面测量~\cite{Aaboud:2016btc}（ID=248）。ATLAS 组完成了共 61 个数据点的 7 项测量：$W^+$（11 点）与 $W^-$（11 点）产生中带电轻子赝快度分布；中央区低质量 Drell-Yan（DY）过程的轻子对快度（6 点）；中央（12 点）与前向（9 点）区的 $Z$ 峰 DY 过程；中央（6 点）与前向（6 点）区的高质量 DY 过程。在已发表的拟合中，我们纳入 3 项测量：$W^+$、$W^-$ 与 $Z$ 峰中央 DY 产生（共 34 点）。
	这些数据只用于拟合 CT18A 与 CT18Z PDF，不用于 CT18 与 CT18X PDF。
	其余数据因可观的电弱修正和/或光子诱导贡献（$\gamma\gamma\to l^+l^-$）而被忽略，讨论见 Sec. \ref{sec:EW}。



	\item 积分亮度 20.3 fb$^{-1}$、$\sqrt{s}=8$ TeV 的 $Z/\gamma^*$ 产生中轻子对横动量 $p_T$ 分布~\cite{Aad:2015auj}（ID=253）。ATLAS 合作组在不变质量范围 $12<M_{\ell \bar \ell}<150$ GeV 内测得轻子对的 $p_{T,  \ell \bar \ell}$ 分布直至 900 GeV。同时实验者给出了单微分分布 $d\sigma/d p_{T,  \ell \bar \ell}$ 与双微分分布 $d^2\sigma/(d p_{T,  \ell \bar \ell}d y_{\ell \bar \ell})$ 的归一化和绝对两种截面。为给 CT18 拟合挑选最干净、最灵敏的数据，
	我们只纳入 $Z$ 峰附近的 3 个不变质量分箱：$M_{\ell \bar \ell}\in[46-66,\ 66-116,\ 116-150]$ GeV。我们不纳入 $M_{\ell \bar \ell} < 46$ GeV 的数据：运动学割断 $p_T^l>20$ GeV 使截面主要来自 $p_{T,  \ell \bar \ell} \gtrsim M_{\ell \bar \ell}$ 的区域，在那里超出当前 ${\cal O}(\alpha_s^3)$ 计算的高阶修正显著。我们既不拟合归一化的 $p_{T,  \ell \bar \ell}$ 分布——因其共享的总归一化会在相距悬殊的 $p_{T,  \ell \bar \ell}$ 区域的粒子计数之间引入人为的相互依赖——也不拟合灵敏度不高的双微分分布。
	我们总共纳入轻子对横动量区域 $45\! \leq\! p_{T,  \ell \bar \ell}\! \leq\! 150$ GeV 内的 27 个数据点，那里固定阶 NNLO 截面最为可靠。更低与更高 $p_{T,  \ell \bar \ell}$ 区域的数据因小 $p_T$ 重求和与电弱修正的贡献而被排除，讨论见 Sec. \ref{sec:EW}。
\end{itemize}

对 CMS，纳入了 $\sqrt{s} = 8$ TeV 单举 ${W}^{\pm}$ 产生电荷不对称的测量~\cite{Khachatryan:2016pev}（ID=249），积分亮度 18.8 fb$^{-1}$。
它们由缪子赝快度的 11 个分箱组成（范围 $0\leq |\eta^\mu|\leq 2.4$），要求 $p^\mu_T \geq 25$ GeV。关联系统误差利用协方差矩阵的分解转换为关联矩阵表示来实现，具体步骤见 Appendix~\ref{sec:chi2_app}。
同样的分解方法也将应用于 LHCb $W/Z$ 实验（Exp ID 245 与 250），把已发表的协方差矩阵转换为关联矩阵。
我们已用既能处理协方差矩阵又能处理关联矩阵表示的 \texttt{ePump}~\cite{Hou:2019gfw} 明确验证了二者的等价性。

在 CT18 中，我们纳入 LHCb 发表的三个实验数据集：
\begin{itemize}
\item $\sqrt{s}=7$ TeV $W/Z$ 前向快度截面测量~\cite{Aaij:2015gna}（ID=245），积分亮度 1.0 fb$^{-1}$，由 $Z$ 玻色子产生截面的 17 个 $Z$ 玻色子快度分箱（$2.0 \leq y_Z \leq 4.25$）与 $W^+$ 或 $W^-$ 玻色子产生的 8 个缪子赝快度分箱（$2.0 \leq \eta^\mu \leq 4.5$）组成。
与上述 CMS 电荷不对称测量类似，系统误差通过把已发表的协方差矩阵转换为关联矩阵来纳入。
束流能量与亮度不确定度在诸截面测量之间视为完全关联。该数据集取代了此前 LHCb 在前向区、积分亮度 35 pb$^{-1}$ 的单举矢量玻色子产生与轻子电荷不对称测量~\cite{Aaij:2012vn}。

\item $\sqrt{s}=8$ TeV 前向快度 $Z\rightarrow e^+e^-$ 截面测量~\cite{Aaij:2015vua}（ID=246），积分亮度 2.0 fb$^{-1}$，由 17 个 $Z$ 玻色子快度分箱（$2.0 \leq y_Z \leq 4.25$）组成。
亮度不确定度取为完全关联，其余关联不确定度则简单按平方相加。我们用 \texttt{ePump} 确认这一近似给出的更新 PDF 与使用协方差矩阵表示所得者相近。


\item $\sqrt{s}=8$ TeV $W/Z$ 产生截面测量~\cite{Aaij:2015zlq}（ID=250），积分亮度 2.0 fb$^{-1}$，由 18 个 $Z$ 玻色子快度分箱（$2.0 \leq y_Z \leq 4.5$）与 $W^+$ 或 $W^-$ 玻色子产生的 8 个缪子赝快度分箱 $(2.0 \leq \eta^\mu \leq 4.5)$ 组成。如同 $\sqrt{s}=7$ TeV 的情形，关联系统误差通过把协方差矩阵转换为关联矩阵表示来纳入。
束流能量与亮度不确定度在诸截面测量之间视为完全关联。

\end{itemize}


\subsubsection{LHC 单举喷注产生
\label{sec:DataJets}
}
对 CMS，使用 $\sqrt{s}=7$ TeV 与 8 TeV 喷注产生的双微分截面测量 $d^2\sigma /(dp_T dy)$。
我们采用两种喷注半径中较大的一个（$R=0.7$）~\cite{Chatrchyan:2014gia}。
数据集包括 $\sqrt{s}=7$ TeV 的 5 fb$^{-1}$（ID=542）与 $\sqrt{s}=8$ TeV 的 19.7 fb$^{-1}$~\cite{Khachatryan:2016mlc}（ID=545）。7 TeV 喷注测量包含六个快度分箱（总快度覆盖 $0\leq |y|\leq 3.0$）中的 158 个数据点，
喷注横动量范围 $56\leq p_T\leq 1327$ GeV。在 8 TeV，同样使用六个快度分箱的喷注数据，覆盖快度范围 $0\leq |y|\leq 3.0$，共 185 个数据点，横动量范围 $74\leq p_T\leq 2500$ GeV。
这些新 CMS 测量取代了此前积分亮度 5 pb$^{-1}$ 的测量~\cite{Chatrchyan:2012bja}。

除 \texttt{HEPData} 文件提供的系统误差信息外，CMS 7 TeV 数据中的喷注能量修正（JEC）已按文献~\cite{Khachatryan:2014waa} 的步骤去关联。特别地，JEC2（``e05''）以及一项额外的、CMS 主张用于 $|y|>2.5$ 的去关联已被实现~\cite{Voutilainen}。这些去关联改善了拟合数据的能力。

对 ATLAS，我们同样采用两种喷注半径中较大的一个（$R=0.6$）。$\sqrt{s}=7$ TeV、$R=0.6$、积分亮度 4.5 fb$^{-1}$ 的单举喷注截面测量~\cite{Aad:2014vwa}（Exp.~ID=544）被纳入全球拟合。
该数据集包含覆盖快度范围 $0.0 \leq y \leq 3.0$ 的六个分箱，在 $74\leq p_T\leq 1992$ GeV 范围共有 140 个数据点。
它取代了此前含 37 pb$^{-1}$ 数据的数据集~\cite{Aad:2011fc}。

按照文献~\cite{Aaboud:2017dvo} 的方案，ATLAS 7 TeV 喷注数据中的两项喷注能量标度（JES）不确定度已去关联，即 MJB（fragmentation）（``jes16''）与味响应（``jes62''）。
去关联步骤使 $\chi^2$ 值降低约 92 个单位。系统误差移动对 $\chi^2$ 的总贡献为 28（对 74 个关联系统误差）。只有一个系统误差源需要大于 2 个标准差的移动。再加入 0.5\% 的理论误差以计及 NNLO 拟合中所需 NNLO/NLO $K$ 因子的蒙特卡洛计算统计噪声~\cite{Currie:2016bfm,Currie:2017ctp,Currie:2018xkj} 后，又获得了 52 个单位的进一步改进，详见 Sec.~\ref{sec:TheorySettings}。
处理 ATLAS 单举喷注数据的更多细节见 Appendix~\ref{sec:ATLASjetdecorrel}。

\subsubsection{LHC 顶夸克对产生}
\label{sec:DataTop}

ATLAS 与 CMS 测量了顶夸克对产生的微分截面，作为顶（反）夸克横向动量 $p_{T,t}$、不变质量 $m_{t\bar{t}}$、正反顶对快度 $y_{t\bar{t}}$、顶夸克对横动量 $p_{T, t\bar{t}}$ 以及顶夸克快度 $y_t$ 的函数；ATLAS 为单独测量，CMS 则成对给出。单个微分 $t\bar{t}$ 截面测量的各自影响已被分析：先用文献~\cite{Wang:2018heo} 的 \texttt{PDFSense} 敏感度框架，再用 \texttt{ePump} 在独立拟合中分析~\cite{Hou:2019gfw,Hou:2019jxd}。各 $t\bar t$ 可观测量之间存在一定张力，导致其对各自偏好的胶子分布有不同的拉动。
同时在 8 TeV 拟合 $p_{T,t}$、$m_{t\bar{t}}$、$y_{t}$ 与 $y_{t\bar{t}}$ 分布的困难也见于~\cite{Bailey:2019yze,Hou:2019gfw,Hou:2019jxd}。

因此在 CT18 分析中，我们决定挑选几个在拟合中相容性最好的顶夸克产生测量。对 ATLAS，利用其发表的统计关联可以纳入不止一个 $t\bar t$ 可观测量。根据 ATLAS 的建议，我们选择了 $\sqrt{s}=8$ TeV、积分亮度 20.3 fb$^{-1}$ 下顶夸克 $p_T$ 的绝对微分截面 $d\sigma/dp_{T,t}$ 与不变质量的 $d\sigma/dm_{t\bar{t}}$（Exp.~ID=580）~\cite{Aad:2015mbv}。\footnote{A. Cooper-Sarkar，私人通信，以及 ATLAS-PHYS-PUB-2018-017。}

这两项 ATLAS 测量合并为一个数据集，包含在 $e$+jets 与 $\mu$+jets 道组合之后、带统计关联的全相空间绝对微分截面（$p_{T,t}$ 与 $m_{t\bar{t}}$ 分布）。这两个分布在把其中一个系统不确定度相对部分子簇射（PS）去关联后一起拟合~\cite{Bogdan}。这些可观测量的 QCD NNLO 理论预言由文献~\cite{Czakon:2017dip,fastnnlo:grids} 提供的 \texttt{fastNNLO} 表获得。


在一项即将开展的研究中我们发现，ATLAS 的单顶夸克与顶夸克对快度分布 $y_t$ 与 $y_{t\bar t}$ 在 CT18 设置下只能以 $\chi^2_E/N_{pt,E}>2.3$ 拟合——过高而不能接受，这与文献~\cite{Kadir:2020yml} 的发现一致。这些分布与其他一些数据集存在张力；无论以单微分还是双微分形式纳入，都不会降低 PDF 不确定度。

对 CMS，我们选择了 $\sqrt{s}=8$ TeV、19.7 fb$^{-1}$ 的归一化双微分截面 $d^2 \sigma/dp_{T,t}dy_t$（Exp.~ID=573）~\cite{Sirunyan:2017azo}。

当 $t\bar t$ 数据集与 Tevatron 及 LHC 喷注产生一起纳入时，它们对 CT18 PDF 的影响是温和的。
其对胶子 PDF 的影响与喷注数据相容，但喷注数据因数据点更多、运动学范围更宽、统计与系统误差相对较小而提供更强的约束。

在 CT18 分析进行期间，CMS 发表了 13 TeV 顶夸克对产生微分截面的测量~\cite{Sirunyan:2018ucr}，且逐分箱数据关联已在 \texttt{HEPData} 库中提供。
虽然这些测量目前未纳入 CT18 全球拟合，对其描述将在 Sec.~\ref{sec:StandardCandles} 讨论。

\subsubsection{未纳入 CT18 拟合的其他 LHC 测量
\label{sec:DataOther}
}

除 CT18(Z) 数据组合之外，我们还仔细研究了 LHC Run-I 的另外几种高亮度测量。在某些情形中，我们要么未观察到显著影响，要么发现与 CT18(Z) 基线存在实质性张力。以下矢量玻色子产生数据曾用 {\tt PDFSense}、{\tt ePump} 或完整拟合作过考察，但未纳入最终的 CT18(Z) 全球分析：
\begin{itemize}
\item 在使理论与 ATLAS $\sqrt{s}=7$ TeV $Z$ 玻色子横向动量分布（$p_{T,  \ell \bar \ell}$）数据（积分亮度 4.7 fb$^{-1}$~\cite{Aad:2014xaa}）良好符合方面遇到困难。这些数据中 $p_{T,  \ell \bar \ell} \sim M_{\ell \bar \ell}$（最适合固定阶计算的运动学区域）的子集在我们试过的各种重正化与因子化标度组合下都给出不可接受的高 $\chi^2$ 值。
在 $Z$ 峰运动学区内，CMS $\sqrt{s}=8$ TeV 的 $p_{T,  \ell \bar \ell}$ 与 $y_{\ell \bar \ell}$ 分布（积分亮度 19.7 fb$^{-1}$~\cite{Khachatryan:2015oaa}）以及 CMS $\sqrt{s}=8$ TeV 归一化 $W$/$Z$ $p_T$ 谱（18.4 pb$^{-1}$~\cite{Khachatryan:2016nbe}）未能赋予显著的约束。
与 CMS 在 $(p_{T,  \ell \bar \ell},y_{\ell \bar \ell})$ 上的双微分分布比较时，我们观察到最后一个快度分箱中理论与数据存在很大分歧。
对于 ATLAS 与 CMS 两组都给出的轻子对{\it 归一化} $p_{T,  \ell \bar \ell}$ 分布，当数据点的归一化依赖拟合范围之外的截面时，如何在有限范围 $p_{T,  \ell \bar \ell} \sim M_{\ell \bar \ell}$ 内与数据一致地比较并不清楚。

\item 无论纳入单微分还是双微分分布 $d \sigma/dQ$ 或 $d^2 \sigma/(dQ\ dy)$，ATLAS $\sqrt{s}=8$ TeV Drell-Yan 截面测量（$116 \leq Q\leq 1500$ GeV、$0 \leq y_Z\leq2.5$，积分亮度 20.3 fb$^{-1}$~\cite{Aad:2016zzw}）都未引起候选 PDF 的实质变化。这些高质量数据受不可忽略的电弱修正与光子诱导（PI）双轻子产生影响，这一点将在 Sec. \ref{sec:EW} 进一步讨论。出于同样的原因，我们不纳入 ATLAS 7 TeV 高质量 Drell-Yan 产生数据（积分亮度 4.7 fb$^{-1}$~\cite{Aad:2013iua}）。



\item ATLAS 合作组 7 TeV 的低质量 Drell-Yan 数据~\cite{Aad:2014qja} 也被考察过，发现对 PDF 无显著影响。

\item 我们考察过 ATLAS~\cite{Aad:2014xca} 与 CMS~\cite{Chatrchyan:2013uja} 在 7 TeV 测量的 $W$ 玻色子伴随粲喷注产生数据的影响。由于 $W+\mbox{charm jet}$ 的 NNLO 计算不可得，这些数据只在 Sec. \ref{sec:Wcharm} 用于与 NLO 理论预言比较。

\item 纳入 CMS 7 TeV~\cite{Chatrchyan:2013tia} 与 8 TeV~\cite{CMS:2014jea} DY 数据的单微分或双微分截面 $d\sigma/dQ$ 或 $d^2\sigma/(dQdy)$ 未发现显著影响。出于如下原因，这些数据未被纳入 CT18 拟合：第一，电弱修正与光子诱导贡献在其高质量区不可忽略；第二，这些数据以全相空间截面形式给出，展开（unfolding）手续带来额外的不确定度。
8 TeV 数据集对 CT18(Z) PDF 有 $\chi^2_E/N_{pt, E} \approx 2$，且经 \texttt{ePump} 与 \texttt{PDFSense} 检验不改变 PDF。
\end{itemize}





\subsection{替代 PDF 拟合：CT18A、CT18X、CT18Z}
\label{sec:alt}
现在我们可以综述与 CT18 并行探索的三个附加拟合，它们对数据选择和理论计算作了不同取舍。这些拟合的关键差异列于表~\ref{tab:AXZ}，其预言将在 Appendix~\ref{sec:AppendixCT18Z} 比较。

\begin{table}
  \begin{tabular}{ccccc}
\hline
\textbf{PDF} & \textbf{Factorization scale } & \textbf{ATLAS 7 TeV $W/Z$} & \textbf{CDHSW $F_{2}^{p,d}$ } & \textbf{Pole charm }\tabularnewline
\textbf{ensemble}
 & \textbf{in DIS} & \textbf{data included?\quad} & \textbf{data included?\quad} & \textbf{mass, GeV}\tabularnewline
\hline
\hline
\noalign{\vskip6pt}
CT18 & $\mu_{F,DIS}^{2}=Q^{2}$ & No & Yes & 1.3\tabularnewline[6pt]
\hline
\noalign{\vskip6pt}
CT18A & $\mu_{F,DIS}^{2}=Q^{2}$ & Yes & Yes & 1.3\tabularnewline[6pt]
\hline
\noalign{\vskip6pt}
CT18X & $\mu_{F,DIS}^{2}=0.8^{2}\left(Q^{2}+\frac{0.3\mbox{ GeV}^{2}}{x_B^{0.3}}\right)$ & No & Yes & 1.3\tabularnewline[6pt]
\hline
\noalign{\vskip6pt}
CT18Z  & $\mu_{F,DIS}^{2}=0.8^{2}\left(Q^{2}+\frac{0.3\mbox{ GeV}^{2}}{x_B^{0.3}}\right)$ & Yes & No & 1.4\tabularnewline[6pt]
\hline
\end{tabular}
        \caption{
                CT18 与三个替代拟合（CT18A、CT18X、CT18Z）的理论设置与数据集选择一览。正文通篇将 CT18Z 与 CT18 比较；关于各替代拟合的更多细节见 App.~\ref{sec:AppendixCT18Z}。
        }
\label{tab:AXZ}
\end{table}


({\it i}) {\bf CT18X} 与 CT18 的差别在于对 DIS 数据集采用了另一种标度选择。最常见的做法是用光子虚度作因子化标度计算单举 DIS 截面，即 $\mu^2_{F,\mathit{DIS}} = Q^2$。然而有论证指出，在 $x\ll 1$ 处对 $\ln^p(1/x)$ 型对数作重求和能把与 HERA Run I+II 数据的符合改善数十个 $\chi^2$ 单位~\cite{Ball:2017otu,Abdolmaleki:2018jln}。在我们的分析中，我们观察到：在固定阶 NNLO 计算中，把 DIS 截面的因子化标度取为调好的 $\mu^2_{F,x} \equiv 0.8^2 \left(Q^2 + 0.3\mbox{ GeV}^2/x_B^{0.3}\right)$ 而非传统的 $\mu^2_F =Q^2$，就能在 HERA DIS 数据集的描述上取得与低 $x$ 重求和分析几乎相同的改进幅度~\cite{Ball:2017otu,Abdolmaleki:2018jln}。以此修改设定所做的拟合记为 CT18X。对该拟合，在 CT18 全球拟合评估的 $Q > 2\mbox{ GeV}$、$x\! >\! 10^{-5}$ 运动学区域内，$\mbox{HERA I+II}$ 的 $\chi^2$ 降低超过 $50$ 个单位。CT18X 对 H1 $F_L$ 的预言适度高于 CT18，使 H1 $F_L$ 的 $\chi^2_E$ 改善几个单位。见图~\ref{fig:saturation} 的示例及其在概要总结~\ref{sec:summary-HERA2} 中的讨论。

依赖 $x_B$ 的标度 $\mu^2_{F,x}$ 的函数形式受饱和图像启发（参见例如
\cite{GolecBiernat:1998js,Caola:2009iy}）。$\mu^2_{F,x}$ 中的数值系数选取为使 HERA DIS 数据的 $\chi^2$ 极小。
在 $x\gtrsim 0.01$，$\mu_{F,x}^2\approx 0.8\ Q^2$ 使 NNLO DIS 截面比取 $\mu_{F}^2=Q^2$ 时更大，这可能模拟次-次-次领头阶（N3LO）及更高阶贡献的效应。在 $x \lesssim 0.01$，$\mu_{F,x}^2$ 数值上减小了 NNLO DIS 截面的 $Q^2$ 导数。这些变化进而导致小 $x$ 处胶子 PDF 增强、$2.5\cdot 10^{-2} \lesssim  x\lesssim 0.2$ 处胶子减弱。

({\it ii}) 与 CT18 不同，{\bf CT18A} 分析纳入了高亮度的 ATLAS 7 TeV $W/Z$ 快度分布~\cite{Aaboud:2016btc}；后者与 DIS 实验有一定张力，并在小 $x_B$ 区偏好比 DIS 实验更大的奇异性 PDF。纳入 ATLAS 7 TeV $W/Z$ 数据使对奇异性 PDF 高度敏感的双缪子 SIDIS 产生数据（NuTeV、CCFR）的 $\chi^2_E$ 值显著恶化（即更大的 $S_E$ 值）。一个察看途径是比较 CT18 拟合的 $S_E$ 分布（图~\ref{fig:sn_ct18}）与 CT18Z 的对应图（App.~\ref{sec:AppendixCT18Z} 图~\ref{fig:sn_ct18z}）。比较表明，由于 CT18Z 拟合纳入了 ATLAS $W/Z$ 数据，CCFR 与 NuTeV 双缪子数据集的 $S_E$ 值相对 CT18 拟合被抬高。另一条验证途径由文献~\cite{Hou:2019gfw} 用 \texttt{ePump} 程序完成。

({\it iii}) {\bf CT18Z} 代表上述设定的累积，旨在获得一个与 CT18 差异最大、同时全局 $\chi^2/N_{pt}$ 与 CT18 大致相当的 PDF 组。CT18Z 拟合同 CT18A 一样纳入 ATLAS 7 TeV $W/Z$ 产生数据，也同 CT18X 一样采用修改的 DIS 标度选择 $\mu_{F,x}$。除此之外，CT18Z 还剔除了来自 $\nu\mathrm{Fe}$ 散射的 CDHSW 结构函数 $F_2$ 与 $x_B F_3$ 提取——否则它们将抵制 CT18Z 在 $x>0.1$ 处较软胶子的趋势，参见 Sec.~\ref{sec:Baseline}。最后，CT18Z 假定略高的粲夸克质量（1.4 GeV，对比 1.3 GeV），以适度改善对矢量玻色子产生数据的拟合。

CT18Z 分析中这些选择的组合使其经由胶子融合的希格斯玻色子产生截面比相应的 CT14 与 CT18 预言低约 1\%。这样，生成四个 CT18(A,X,Z) 拟合过程中所做的种种选择，使我们能更忠实地探索与现有强子数据相容的 NNLO PDF 行为的全部范围，并对电弱精密物理具有意义。
